\begin{statementbox}
\begin{center}\begin{tikzpicture}[font=\small,>=stealth,line join=round]
\draw[gray!40,rounded corners=3pt] (0,0) rectangle (14.2,8.8);\node at (7.1,8.200000000000001) {奇配轴、偶配点：隔2变号，隔4回到原高度};
\node at (3.1500000000000004,7.3) {奇函数配轴 $x=1$};\begin{scope}[shift={(2.061111111111111,4.345714285714285)},x=0.5444444444444445cm,y=1.0857142857142856cm]\draw[->,gray] (-2.5,0)--(6.5,0) node[right]{$x$};\draw[->,gray] (0,-1.7)--(0,1.8) node[above]{$y$};\draw[blue,thick,solid,domain=-2:-1,samples=180] plot(\x,{-1*\x+(-2)});\draw[blue,thick,solid,domain=-1:0,samples=180] plot(\x,{1*\x+(0)});\draw[blue,thick,solid,domain=0:1,samples=180] plot(\x,{1*\x+(0)});\draw[blue,thick,solid,domain=1:2,samples=180] plot(\x,{-1*\x+(2)});\draw[blue,thick,solid,domain=2:3,samples=180] plot(\x,{-1*\x+(2)});\draw[blue,thick,solid,domain=3:4,samples=180] plot(\x,{1*\x+(-4)});\draw[blue,thick,solid,domain=4:5,samples=180] plot(\x,{1*\x+(-4)});\draw[blue,thick,solid,domain=5:6,samples=180] plot(\x,{-1*\x+(6)});\fill[orange!85!black] (1,1) circle(2pt);\fill[orange!85!black] (3,-1) circle(2pt);\fill[orange!85!black] (5,1) circle(2pt);\draw[violet,dashed] (1,-1.4)--(1,1.4);\end{scope}\node[align=center,text width=5.8cm] at (3.1500000000000004,1.7) {输入 $1,3,5$ 对应 $1,-1,1$\\一个正周期为4};\node at (9.95,7.3) {偶函数配中心 $(1,0)$};\begin{scope}[shift={(8.86111111111111,4.345714285714285)},x=0.5444444444444445cm,y=1.0857142857142856cm]\draw[->,gray] (-2.5,0)--(6.5,0) node[right]{$x$};\draw[->,gray] (0,-1.7)--(0,1.8) node[above]{$y$};\draw[teal,thick,solid,domain=-2:-1,samples=180] plot(\x,{1*\x+(1)});\draw[teal,thick,solid,domain=-1:0,samples=180] plot(\x,{1*\x+(1)});\draw[teal,thick,solid,domain=0:1,samples=180] plot(\x,{-1*\x+(1)});\draw[teal,thick,solid,domain=1:2,samples=180] plot(\x,{-1*\x+(1)});\draw[teal,thick,solid,domain=2:3,samples=180] plot(\x,{1*\x+(-3)});\draw[teal,thick,solid,domain=3:4,samples=180] plot(\x,{1*\x+(-3)});\draw[teal,thick,solid,domain=4:5,samples=180] plot(\x,{-1*\x+(5)});\draw[teal,thick,solid,domain=5:6,samples=180] plot(\x,{-1*\x+(5)});\fill[orange!85!black] (0,1) circle(2pt);\fill[orange!85!black] (2,-1) circle(2pt);\fill[orange!85!black] (4,1) circle(2pt);\fill[orange!85!black] (1,0) circle(2pt);\end{scope}\node[align=center,text width=5.8cm] at (9.95,1.7) {输入 $0,2,4$ 对应 $1,-1,1$\\一个正周期为4};
\node[text width=12.6cm,align=center] at (7.1,.65) {两图同尺度，均为真实分段直线函数；高低相反时还不能叫周期。};
\end{tikzpicture}\end{center}


\textbf{\textcolor{CStatement}{这条结论解决什么问题？}}

有些题目没有告诉我们 $f(x)$ 的具体式子，只说它是奇函数、偶函数，或者图像有一条对称轴、一个对称中心。怎样在不知道函数式子的情况下，判断函数会不会重复出现？这条结论就是把\textbf{两条对称信息接起来，找出重复的距离}。

\medskip
\end{statementbox}\clearpage\begin{statementbox}\textbf{\textcolor{CStatement}{先记住两个搭配}}（以下先假设函数定义在整个实数范围 $\mathbb{R}$ 上，$a\ne0$）
\begin{enumerate}[leftmargin=*,itemsep=0.3em]
  \item \textbf{奇函数 $+$ 关于直线 $x=a$ 对称}：
  \[f(-x)=-f(x),\qquad f(a+t)=f(a-t).\]
  \item \textbf{偶函数 $+$ 关于点 $(a,0)$ 对称}：
  \[f(-x)=f(x),\qquad f(a+t)=-f(a-t).\]
\end{enumerate}

\textbf{\textcolor{CStatement}{这两种情况，都能推出同一个结论}}
\[
\boxed{f(x+2a)=-f(x)}
\qquad\Longrightarrow\qquad
\boxed{f(x+4a)=f(x)}.
\]

\textbf{通俗地说：}沿 $x$ 轴移动 $2|a|$，函数值变成原来的相反数；再移动 $2|a|$，函数值恢复原样。因此
\[
\boxed{4|a|\text{ 一定是函数的一个正周期。}}
\]
当 $a<0$ 时，上式最初得到的是向左移动；改写为向右移动 $2|a|$ 和 $4|a|$ 仍然成立。

\medskip
\end{statementbox}\clearpage\begin{statementbox}\textbf{\textcolor{CStatement}{特别重要：不要把“一个周期”写成“最小正周期”}}

仅凭上述条件，\textbf{无法确定最小正周期一定是 $4|a|$}。函数可能在更短的距离就重复；零函数甚至没有最小正周期。求函数值时，只要知道 $4|a|$ 是一个周期，就已经够用了。

\medskip
\textbf{\textcolor{CStatement}{适用边界}}
\begin{itemize}[leftmargin=*,itemsep=0.25em]
  \item 这里取定义域为 $\mathbb{R}$，便于把每一次代入都用于任意实数。对于其他定义域，须另外保证每次对应的点都在定义域内，而且整体平移后仍有意义。
  \item $a=0$ 不能当成一个正的距离：在这两种搭配下会推出 $f(x)\equiv0$，不能由 $4|a|=0$ 得出正周期。
\end{itemize}
\end{statementbox}
